\documentclass[10pt,a4paper]{amsart}
\usepackage{amsmath,amssymb,amsthm,graphicx}
\usepackage{verbatim}
\newenvironment{reference}{\comment}{\endcomment}
%\newenvironment{reference}{}{}
\newenvironment{slogan}{\comment}{\endcomment}
\newenvironment{history}{\comment}{\endcomment}

% For commutative diagrams we use Xy-pic
\usepackage[all]{xy}

% We use 2cell for 2-commutative diagrams.
\xyoption{2cell}
% Theorem environments.
%
\theoremstyle{plain}
\newtheorem{theorem}[subsection]{Theorem}
\newtheorem{proposition}[subsection]{Proposition}
\newtheorem{lemma}[subsection]{Lemma}

\theoremstyle{definition}
\newtheorem{definition}[subsection]{Definition}
\newtheorem{example}[subsection]{Example}
\newtheorem{exercise}[subsection]{Exercise}
\newtheorem{situation}[subsection]{Situation}

\theoremstyle{remark}
\newtheorem{remark}[subsection]{Remark}
\newtheorem{remarks}[subsection]{Remarks}

\numberwithin{equation}{subsection}

% Macros
%
\def\lim{\mathop{\mathrm{lim}}\nolimits}
\def\colim{\mathop{\mathrm{colim}}\nolimits}
\def\Spec{\mathop{\mathrm{Spec}}}
\def\Hom{\mathop{\mathrm{Hom}}\nolimits}
\def\Ext{\mathop{\mathrm{Ext}}\nolimits}
\def\SheafHom{\mathop{\mathcal{H}\!\mathit{om}}\nolimits}
\def\SheafExt{\mathop{\mathcal{E}\!\mathit{xt}}\nolimits}
\def\Sch{\mathit{Sch}}
\def\Mor{\mathop{\mathrm{Mor}}\nolimits}
\def\Ob{\mathop{\mathrm{Ob}}\nolimits}
\def\Sh{\mathop{\mathit{Sh}}\nolimits}
\def\NL{\mathop{N\!L}\nolimits}
\def\CH{\mathop{\mathrm{CH}}\nolimits}
\def\proetale{{pro\text{-}\acute{e}tale}}
\def\etale{{\acute{e}tale}}
\def\QCoh{\mathit{QCoh}}
\def\Ker{\mathop{\mathrm{Ker}}}
\def\Im{\mathop{\mathrm{Im}}}
\def\Coker{\mathop{\mathrm{Coker}}}
\def\Coim{\mathop{\mathrm{Coim}}}

% Boxtimes
%
\DeclareMathSymbol{\boxtimes}{\mathbin}{AMSa}{"02}

%
% Macros for moduli stacks/spaces
%
\def\QCohstack{\mathcal{QC}\!\mathit{oh}}
\def\Cohstack{\mathcal{C}\!\mathit{oh}}
\def\Spacesstack{\mathcal{S}\!\mathit{paces}}
\def\Quotfunctor{\mathrm{Quot}}
\def\Hilbfunctor{\mathrm{Hilb}}
\def\Curvesstack{\mathcal{C}\!\mathit{urves}}
\def\Polarizedstack{\mathcal{P}\!\mathit{olarized}}
\def\Complexesstack{\mathcal{C}\!\mathit{omplexes}}
% \Pic is the operator that assigns to X its picard group, usage \Pic(X)
% \Picardstack_{X/B} denotes the Picard stack of X over B
% \Picardfunctor_{X/B} denotes the Picard functor of X over B
\def\Pic{\mathop{\mathrm{Pic}}\nolimits}
\def\Picardstack{\mathcal{P}\!\mathit{ic}}
\def\Picardfunctor{\mathrm{Pic}}
\def\Deformationcategory{\mathcal{D}\!\mathit{ef}}
\usepackage{cleveref}
\crefname{originallemma}{original lemma}{original lemmas}
\Crefname{originallemma}{Original lemma}{Original lemmas}
\makeatletter
\newlabel{coherent-lemma-proper-over-affine-cohomology-finite}{{02O6}{source}{}{sourcepointer-1}{}}
\newlabel{coherent-lemma-proper-over-affine-cohomology-finite@cref}{{[originallemma][0][]02O6}{[1][0][]source}{}{}{}}
\makeatother
\crefname{originalsection}{original section}{original sections}
\Crefname{originalsection}{Original section}{Original sections}
\makeatletter
\newlabel{duality-example-affine-twisted-inverse-image}{{0A9F}{source}{}{sourcepointer-2}{}}
\newlabel{duality-example-affine-twisted-inverse-image@cref}{{[originalsection][0][]0A9F}{[1][0][]source}{}{}{}}
\makeatother
\crefname{originallemma}{original lemma}{original lemmas}
\Crefname{originallemma}{Original lemma}{Original lemmas}
\makeatletter
\newlabel{duality-lemma-more-base-change}{{0AA8}{source}{}{sourcepointer-3}{}}
\newlabel{duality-lemma-more-base-change@cref}{{[originallemma][0][]0AA8}{[1][0][]source}{}{}{}}
\makeatother
\crefname{originallemma}{original lemma}{original lemmas}
\Crefname{originallemma}{Original lemma}{Original lemmas}
\makeatletter
\newlabel{duality-lemma-twisted-inverse-image}{{0A9E}{source}{}{sourcepointer-4}{}}
\newlabel{duality-lemma-twisted-inverse-image@cref}{{[originallemma][0][]0A9E}{[1][0][]source}{}{}{}}
\makeatother
\crefname{originalremark}{original remark}{original remarks}
\Crefname{originalremark}{Original remark}{Original remarks}
\makeatletter
\newlabel{duality-remark-relative-dualizing-complex}{{0B6S}{source}{}{sourcepointer-5}{}}
\newlabel{duality-remark-relative-dualizing-complex@cref}{{[originalremark][0][]0B6S}{[1][0][]source}{}{}{}}
\makeatother
\crefname{originalreference}{original reference}{original references}
\Crefname{originalreference}{Original reference}{Original references}
\makeatletter
\newlabel{lemma-duality-dim-1}{{0BS2}{source}{}{sourcepointer-6}{}}
\newlabel{lemma-duality-dim-1@cref}{{[originalreference][0][]0BS2}{[1][0][]source}{}{}{}}
\makeatother
\title{The dualizing complex and sheaf of a proper curve behave correctly under intermediate fields and base change}
\author{The Stacks Project Authors}
\date{stacks-project:curves.tex:local-snapshot-20261012; GNU FDL 1.2 or later}
\begin{document}
\maketitle
\noindent\textit{Editorial note.} All delivered mathematical text is the author's own, copied from the pinned Stacks Project chapter file curves.tex; no statement or proof step was rewritten, reordered or removed.\par\medskip
\noindent\textit{Editorial note.} The counted proof is the author's complete argument: the adjoint reformulation, the finiteness input, the uniqueness-of-adjoints decomposition, the Hom identification by dimension count, and the base-change derivation of part (2) are all retained verbatim.\par\medskip
\noindent\textit{Editorial note.} The label lemma-sanity-check-duality names the counted statement, and the proof environment immediately following it in the source is delivered verbatim as its proof.\par\medskip
\noindent\textit{Editorial note.} The counted statement is carried by the author's own lemma environment in the chapter source.\par\medskip
\noindent\textit{Allowed standard prerequisites (editorial).}\begin{itemize}
\item The duality framework being tested. Same chapter, Lemma lemma-duality-dim-1. Fixes omega\_X\textasciicircum{}bullet = a(O\_Spec(k)) with a the right adjoint of the twisted inverse image functor, the object whose properties are rechecked.
\item Twisted inverse images and their adjoints. Duality for Schemes, Lemma duality-lemma-twisted-inverse-image, Example duality-example-affine-twisted-inverse-image, Lemma duality-lemma-more-base-change and Remark duality-remark-relative-dualizing-complex. Supplies the adjoint decompositions f\textasciicircum{}! = (f')\textasciicircum{}! composed with g\textasciicircum{}!, the affine computation Hom\_k(k', k) = k', and the base change theorem for a.
\item Finiteness of global functions. Cohomology of Schemes, Lemma coherent-lemma-proper-over-affine-cohomology-finite. Shows k'/k is finite since k' sits inside H\textasciicircum{}0(X, O\_X) of the proper scheme X.
\end{itemize}\medskip
\ifdefined\hypertarget\hypertarget{sourcepointer-1}{}\fi
\noindent\textit{Source reference.} Original lemma 02O6: sources/stacks-chow/tags.txt, source lines 3505--3505. Prerequisite pointer invoked at proof line(s) 957; it names a standard fact of a sibling chapter that the counted argument uses, and every step of the counted proof built on it is retained verbatim in the delivered body.\par\medskip
\ifdefined\hypertarget\hypertarget{sourcepointer-2}{}\fi
\noindent\textit{Source reference.} Original section 0A9F: sources/stacks-chow/tags.txt, source lines 12948--12948. Prerequisite pointer invoked at proof line(s) 964; it names a standard fact of a sibling chapter that the counted argument uses, and every step of the counted proof built on it is retained verbatim in the delivered body.\par\medskip
\ifdefined\hypertarget\hypertarget{sourcepointer-3}{}\fi
\noindent\textit{Source reference.} Original lemma 0AA8: sources/stacks-chow/tags.txt, source lines 12976--12976. Prerequisite pointer invoked at proof line(s) 970; it names a standard fact of a sibling chapter that the counted argument uses, and every step of the counted proof built on it is retained verbatim in the delivered body.\par\medskip
\ifdefined\hypertarget\hypertarget{sourcepointer-4}{}\fi
\noindent\textit{Source reference.} Original lemma 0A9E: sources/stacks-chow/tags.txt, source lines 12947--12947. Prerequisite pointer invoked at proof line(s) 949; it names a standard fact of a sibling chapter that the counted argument uses, and every step of the counted proof built on it is retained verbatim in the delivered body.\par\medskip
\ifdefined\hypertarget\hypertarget{sourcepointer-5}{}\fi
\noindent\textit{Source reference.} Original remark 0B6S: sources/stacks-chow/tags.txt, source lines 14080--14080. Prerequisite pointer invoked at proof line(s) 972; it names a standard fact of a sibling chapter that the counted argument uses, and every step of the counted proof built on it is retained verbatim in the delivered body.\par\medskip
\ifdefined\hypertarget\hypertarget{sourcepointer-6}{}\fi
\noindent\textit{Source reference.} Original reference 0BS2: sources/stacks-curves/pinned-original-source.tex, source lines 796--796. Prerequisite pointer invoked at proof line(s) 929; it names a standard fact of the same chapter that the counted argument uses, and every step of the counted proof built on it is retained verbatim in the delivered body.\par\medskip
\begin{lemma}
\label{lemma-sanity-check-duality}
Let $X$ be a proper scheme of dimension $\leq 1$ over a field $k$.
Let $\omega_X^\bullet$ and $\omega_X$ be as in Lemma \ref{lemma-duality-dim-1}.
\begin{enumerate}
\item If $X \to \Spec(k)$ factors as $X \to \Spec(k') \to \Spec(k)$
for some field $k'$, then $\omega_X^\bullet$ and $\omega_X$
satisfy properties (4), (5), (6) with $k$ replaced with $k'$.
\item If $K/k$ is a field extension, then the pullback of
$\omega_X^\bullet$ and $\omega_X$ to the base change $X_K$
are as in  Lemma \ref{lemma-duality-dim-1} for the morphism
$X_K \to \Spec(K)$.
\end{enumerate}
\end{lemma}

\begin{proof}
Denote $f : X \to \Spec(k)$ the structure morphism.
Assertion (1) really means that $\omega_X^\bullet$ and $\omega_X$
are as in Lemma \ref{lemma-duality-dim-1} for the morphism
$f' : X \to \Spec(k')$. In the proof of Lemma \ref{lemma-duality-dim-1}
we took $\omega_X^\bullet = a(\mathcal{O}_{\Spec(k)})$
where $a$ be is the right adjoint of
Duality for Schemes, Lemma
\ref{duality-lemma-twisted-inverse-image} for $f$.
Thus we have to show
$a(\mathcal{O}_{\Spec(k)}) \cong a'(\mathcal{O}_{\Spec(k)})$
where $a'$ be is the right adjoint of
Duality for Schemes, Lemma
\ref{duality-lemma-twisted-inverse-image} for $f'$.
Since $k' \subset H^0(X, \mathcal{O}_X)$ we see that $k'/k$ is a finite
extension (Cohomology of Schemes, Lemma
\ref{coherent-lemma-proper-over-affine-cohomology-finite}).
By uniqueness of adjoints we have $a = a' \circ b$ where
$b$ is the right adjoint of Duality for Schemes, Lemma
\ref{duality-lemma-twisted-inverse-image} for $g : \Spec(k') \to \Spec(k)$.
Another way to say this: we have $f^! = (f')^! \circ g^!$.
Thus it suffices to show that $\Hom_k(k', k) \cong k'$ as
$k'$-modules, see Duality for Schemes, Example
\ref{duality-example-affine-twisted-inverse-image}.
This holds because these are $k'$-vector spaces of
the same dimension (namely dimension $1$).

\medskip\noindent
Proof of (2). This holds because we have base change for $a$ by
Duality for Schemes, Lemma \ref{duality-lemma-more-base-change}.
See discussion in Duality for Schemes, Remark
\ref{duality-remark-relative-dualizing-complex}.
\end{proof}

\end{document}
